\documentclass[11pt]{article}

\usepackage{amsmath}
\usepackage{amssymb}
\usepackage{amsthm}
\usepackage{mathtools}
\usepackage{stmaryrd}
\usepackage[margin=1in]{geometry}
\usepackage{tikz-cd}
\usepackage{hyperref}

% ---- Theorem environments ----
\theoremstyle{plain}
\newtheorem{theorem}{Theorem}[section]
\newtheorem{lemma}[theorem]{Lemma}
\newtheorem{proposition}[theorem]{Proposition}
\newtheorem{corollary}[theorem]{Corollary}

\theoremstyle{definition}
\newtheorem{definition}[theorem]{Definition}
\newtheorem{example}[theorem]{Example}
\newtheorem{remark}[theorem]{Remark}
\newtheorem{notation}[theorem]{Notation}

% ---- Notation macros ----
\newcommand{\Set}{\mathbf{Set}}
\newcommand{\Cat}{\mathbf{Cat}}
\newcommand{\Cont}{\mathbf{Cont}}
\newcommand{\Fam}{\mathbf{Fam}}
\newcommand{\op}{^{\mathrm{op}}}
\newcommand{\yA}{y^{A}}
\newcommand{\Nat}{\mathrm{Nat}}
\newcommand{\Lan}{\mathrm{Lan}}
\newcommand{\Ran}{\mathrm{Ran}}
\newcommand{\id}{\mathrm{id}}
\newcommand{\ob}{\mathrm{ob}}
\newcommand{\co}{\mathbin{\text{\Large$\cdot$}}}      % copower dot
\newcommand{\Day}{\circledast}
\newcommand{\Dayc}{\odot}
\newcommand{\tri}{\triangleleft}       % sequential (composition) operator, cf. monoidal chapter
\DeclareMathOperator{\Hom}{Hom}
\newcommand{\iso}{\cong}
\newcommand{\ext}[1]{\llbracket #1 \rrbracket}

\title{Chapter 0 --- The Machinery:\\ Representables, Yoneda, Day Convolution, and Kan Extensions}
\author{MacBeth}
\date{}

% Book conventions: this is a chapter of \emph{The Category of Containers}.
% Cross-references of the form ``Theorem~C'' point to the four-monoidal chapter.

\begin{document}
\maketitle

\begin{abstract}
This chapter assembles the pre-container tools the rest of the book leans on:
representable functors and the Yoneda lemma, density of representables, left
Kan extensions and the coend calculus, Day convolution, and the interaction of
the Yoneda embedding with monoidal and closed structure. \emph{Everything here
is standard}, collected for self-containedness and to fix notation. We claim no
originality: the Yoneda lemma is Yoneda's and Mac Lane's; Kan extensions and the
coend calculus follow Mac Lane~\cite{MacLane} Ch.~X, Riehl~\cite{Riehl}, and
Loregian~\cite{Loregian}; Day convolution is Day's~\cite{Day1970}. Our only
contribution is arrangement --- pedagogical order, and proofs written in our own
words --- and the deliberate foregrounding of two observations of Neil Ghani that
the book's monoidal chapter will invoke as black boxes: that free coproduct
completion \emph{is} a left Kan extension along the representable embedding
(\S\ref{sec:kan}), and that closed structure on a Day-convolution category is
\emph{forced on representables and extended by density}
(\S\ref{sec:closed}). We are careful about the second: density determines the
internal hom, but it does \emph{not} follow that an arbitrary cocontinuous
strong-monoidal functor preserves it (we give a one-line counterexample,
Remark~\ref{rem:notfree}); preservation is a genuine, separate fact, true for the
container-extension functor and checked directly in the monoidal chapter.
\end{abstract}

\tableofcontents

\bigskip
\noindent\textbf{Conventions.} $\Set$ is the category of sets and functions.
For a locally small category $\mathcal{C}$ we write $\mathcal{C}(a,b)$ for its
hom-set. All functors are covariant unless a superscript $\op$ says otherwise. A
\emph{presheaf} on $\mathcal{C}$ is a functor $\mathcal{C}\op\to\Set$; a
\emph{copresheaf} is a functor $\mathcal{C}\to\Set$. We use $\iso$ for
canonical (natural) isomorphism throughout, and we are careful to say
\emph{which} isomorphism whenever it matters.

% =====================================================================
\section{Representable functors and the Yoneda lemma}
\label{sec:yoneda}
% =====================================================================

Every construction in this book eventually reduces to a hom-set. The
representable functors are the functors that \emph{are} a hom-set, one argument
held fixed, and they turn out to generate everything else. We start there.

\begin{definition}[Representable copresheaf]
\label{def:rep}
Fix a set $A$. The \emph{covariant representable} at $A$ is the functor
\[
  \yA \;:=\; \Set(A,-)\;:\;\Set \longrightarrow \Set,
  \qquad
  \yA(X) = \Set(A,X) = X^{A},
\]
sending a function $f\colon X\to Y$ to post-composition
$\yA(f) = f\circ(-)\colon X^{A}\to Y^{A}$. A copresheaf $F\colon\Set\to\Set$ is
\emph{representable} if $F\iso\yA$ for some $A$; we call $A$ the
\emph{representing object}.
\end{definition}

\begin{remark}[The container reading]
\label{rem:containerrep}
This book is about \emph{containers}. A container is a pair $(S,P)$ of a set $S$
of \emph{shapes} and a family $P\colon S\to\Set$ of \emph{position sets}, and its
\emph{extension} is the functor
\[
  \ext{S,P}\;:\;\Set\to\Set,
  \qquad
  \ext{S,P}(X) \;=\; \sum_{s\in S}\Set\!\bigl(P(s),X\bigr).
\]
The \emph{representable container} is the one with a single shape and $A$
positions, $(S,P)=(1,\,A)$ (the family $1\to\Set$ picking out $A$). Its extension
is
\[
  \ext{1,A}(X) \;=\; \sum_{s\in 1}\Set(A,X) \;=\; \Set(A,X) \;=\; X^{A} \;=\; \yA(X).
\]
So representable copresheaves are exactly the extensions of single-shape
containers, and the coproduct in the general formula is a coproduct \emph{of
representables} indexed by shapes. This is not a coincidence to be admired and
forgotten: \S\ref{sec:kan} and \S\ref{sec:day} show that ``a container is a
coproduct of representables'' is the whole story, phrased first as a Kan
extension and then as a Day convolution.
\end{remark}

\subsection{The Yoneda embedding and the Yoneda lemma}

Definition~\ref{def:rep} is uniform in $A$: a function $a\colon A'\to A$ induces,
by \emph{pre}-composition, a natural transformation $\yA\Rightarrow y^{A'}$. Thus
$A\mapsto\yA$ is a functor $\Set\op\to[\Set,\Set]$, the \emph{Yoneda embedding}
(covariant version). We write it $y^{(-)}$. It is genuinely an embedding: fully
faithful, as the Yoneda lemma will show.

\begin{theorem}[Yoneda lemma]
\label{thm:yoneda}
Let $F\colon\Set\to\Set$ be any functor and $A$ any set. There is a bijection
\[
  \Theta_{F,A}\;:\;\Nat(\yA,\,F)\;\xrightarrow{\ \iso\ }\;F(A),
\]
natural in both $F$ and $A$, given by evaluating a natural transformation at the
identity: $\Theta(\alpha) = \alpha_{A}(\id_{A})$.
\end{theorem}

\begin{proof}
We build the inverse and check both round trips.

\emph{Inverse.} Given an element $x\in F(A)$, define a family of maps
$\Psi(x)_{X}\colon \yA(X)=\Set(A,X)\to F(X)$ by
\[
  \Psi(x)_{X}(g) \;=\; F(g)(x),\qquad g\colon A\to X .
\]
This is natural in $X$: for $h\colon X\to Y$ we must show
$F(h)\circ\Psi(x)_{X} = \Psi(x)_{Y}\circ\yA(h)$. Applied to $g\colon A\to X$ the
left side is $F(h)(F(g)(x)) = F(h\circ g)(x)$ by functoriality, and the right side
is $\Psi(x)_{Y}(h\circ g) = F(h\circ g)(x)$. They agree, so $\Psi(x)$ is a
natural transformation $\yA\Rightarrow F$.

\emph{$\Theta$ after $\Psi$.} For $x\in F(A)$,
\[
  \Theta(\Psi(x)) = \Psi(x)_{A}(\id_{A}) = F(\id_{A})(x) = \id_{F(A)}(x) = x .
\]

\emph{$\Psi$ after $\Theta$.} Let $\alpha\colon\yA\Rightarrow F$ and put
$x=\Theta(\alpha)=\alpha_{A}(\id_{A})$. We must show $\Psi(x)=\alpha$, i.e.\ for
every $X$ and every $g\colon A\to X$ that $F(g)(x) = \alpha_{X}(g)$. Naturality of
$\alpha$ at the map $g\colon A\to X$ gives a commuting square
\[
\begin{tikzcd}[column sep=large]
  \Set(A,A) \arrow[r,"\alpha_{A}"] \arrow[d,"g\circ(-)"'] & F(A) \arrow[d,"F(g)"]\\
  \Set(A,X) \arrow[r,"\alpha_{X}"'] & F(X).
\end{tikzcd}
\]
Chase $\id_{A}$ around it. Down-then-right gives
$\alpha_{X}(g\circ\id_{A}) = \alpha_{X}(g)$; right-then-down gives
$F(g)(\alpha_{A}(\id_{A})) = F(g)(x)$. Hence $\alpha_{X}(g) = F(g)(x) =
\Psi(x)_{X}(g)$, as required.

Thus $\Theta$ and $\Psi$ are mutually inverse bijections. Naturality of $\Theta$
in $F$ (along natural transformations $F\Rightarrow F'$) and in $A$ (along
functions $A'\to A$, which act on $\yA$ by pre-composition) is a routine diagram
chase of the same shape; we omit it. \end{proof}

\begin{corollary}[Full faithfulness]
\label{cor:ff}
Taking $F=y^{B}$ in Theorem~\ref{thm:yoneda} gives
$\Nat(\yA,y^{B})\iso y^{B}(A)=\Set(B,A)$, natural in $A$ and $B$. Hence the
Yoneda embedding $y^{(-)}\colon\Set\op\to[\Set,\Set]$ is fully faithful: natural
transformations between representables are exactly (pre-composition by)
functions, in the opposite direction.
\end{corollary}

\begin{example}[Yoneda in $\Set$, concretely]
\label{ex:yoneda}
Take $A=1$, the one-element set. Then $\yA = \Set(1,-)$ is naturally isomorphic to
the identity functor, since $\Set(1,X)\iso X$. The Yoneda lemma reads
$\Nat(\mathrm{Id},F)\iso F(1)$: a natural self-map of the identity is pinned down
by a single element of $F(1)$. For $F=\mathrm{Id}$ this recovers
$\Nat(\mathrm{Id},\mathrm{Id})\iso 1$: the only natural transformation
$\mathrm{Id}\Rightarrow\mathrm{Id}$ is the identity. Try to build a
``doubling'' natural transformation $X\to X$ and you cannot even start: there is
no natural way to duplicate an anonymous element, precisely because $\id_{1}$ is
the only probe you are allowed and it must go to itself.
\end{example}

\subsection{Density of representables}
\label{sub:density}

The Yoneda lemma says representables see everything. Density says representables
\emph{are} everything, once you are allowed colimits.

\begin{theorem}[Density / co-Yoneda]
\label{thm:density}
Every functor $F\colon\Set\to\Set$ is canonically a colimit of representables.
Concretely, $F$ is the coend
\[
  F \;\iso\; \int^{A\in\Set} F(A)\cdot \yA,
\]
where $F(A)\cdot\yA$ denotes the $F(A)$-fold copower (coproduct) of the
representable $\yA$. Equivalently, $F$ is the colimit of the canonical diagram
\[
  \bigl(\, \yA \;\xrightarrow{\ \Psi(x)\ }\; F \,\bigr)_{(A,\,x),\ x\in F(A)}
\]
indexed by the category of elements of $F$; the coend presentation packages this
diagram and its identifications.
\end{theorem}

\begin{proof}[Proof sketch]
This is the ``density formula,'' dual to Yoneda; see
Riehl~\cite[\S6.2]{Riehl} or Loregian~\cite[Ch.~2]{Loregian}. The universal
property is exactly Theorem~\ref{thm:yoneda} read backwards: a cocone from the
diagram $F(A)\cdot\yA$ to a functor $G$ is, by the copower and Yoneda
adjunctions, a family of maps $F(A)\to \Nat(\yA,G)\iso G(A)$ natural in $A$, i.e.\
exactly a natural transformation $F\Rightarrow G$. Since cocones out of the coend
biject naturally with maps $F\Rightarrow G$, the coend \emph{is} $F$. The coend's
coequalizer presentation
\[
  \coprod_{f\colon A\to A'} F(A)\cdot y^{A'}
  \;\rightrightarrows\;
  \coprod_{A} F(A)\cdot \yA
  \;\twoheadrightarrow\; F
\]
records the identification ``push $x$ along $f$ in $F$'' $=$ ``reindex the
representable along $f$,'' which is where naturality is spent. \end{proof}

\begin{remark}
\label{rem:densegen}
Density is the load-bearing fact for \S\ref{sec:kan} and \S\ref{sec:closed}: the
representables form a \emph{dense generator} of $[\Set,\Set]$. Any two of our
constructions that agree on representables and both preserve the relevant
colimits agree everywhere. We will use this uniformity of argument for Neil's
observation~(ii) --- ``the internal hom is determined on representables and
extended by density.''
\end{remark}

% =====================================================================
\section{Left Kan extensions}
\label{sec:kan}
% =====================================================================

A Kan extension is the best approximation to ``extend a functor along another
functor.'' We need the left version and its pointwise coend formula, and then one
theorem: for our purposes the free coproduct completion of a category \emph{is} a
left Kan extension along the representable embedding. This is the lemma the
monoidal chapter cites for its Theorem~C.

\begin{definition}[Left Kan extension]
\label{def:lan}
Let $K\colon\mathcal{C}\to\mathcal{D}$ and $F\colon\mathcal{C}\to\mathcal{E}$ be
functors. A \emph{left Kan extension} of $F$ along $K$ is a functor
$\Lan_{K}F\colon\mathcal{D}\to\mathcal{E}$ together with a natural transformation
$\eta\colon F\Rightarrow (\Lan_{K}F)\circ K$ that is universal from $F$: for every
$G\colon\mathcal{D}\to\mathcal{E}$ and every $\gamma\colon F\Rightarrow G\circ K$
there is a \emph{unique} $\bar\gamma\colon\Lan_{K}F\Rightarrow G$ with
$(\bar\gamma K)\circ\eta = \gamma$.
\[
\begin{tikzcd}[column sep=large,row sep=large]
  \mathcal{C} \arrow[r,"F"] \arrow[d,"K"'] & \mathcal{E}\\
  \mathcal{D} \arrow[ur,"\Lan_{K}F"',""{name=L,above}] &
  \arrow[from=1-1,to=2-1,phantom]
  \arrow[Rightarrow,from=1-2,to=L,"\eta"',shorten <=6pt,shorten >=6pt]
\end{tikzcd}
\]
Equivalently, $\Lan_{K}(-)$ is left adjoint to precomposition
$(-)\circ K\colon[\mathcal{D},\mathcal{E}]\to[\mathcal{C},\mathcal{E}]$.
\end{definition}

\begin{theorem}[Pointwise coend formula]
\label{thm:lancoend}
If $\mathcal{C}$ is small, $\mathcal{E}$ is cocomplete, and $\mathcal{D}$ is
locally small, then $\Lan_{K}F$ exists and is computed pointwise by the coend
\[
  (\Lan_{K}F)(d) \;\iso\; \int^{c\in\mathcal{C}} \mathcal{D}(Kc,\,d)\cdot F(c),
\]
where $\mathcal{D}(Kc,d)\cdot F(c)$ is the $\mathcal{D}(Kc,d)$-fold copower of the
object $F(c)\in\mathcal{E}$.
\end{theorem}

\begin{proof}[Reference]
Standard; Mac Lane~\cite[Ch.~X, \S4]{MacLane}, Riehl~\cite[\S6.2]{Riehl},
Loregian~\cite[Ch.~2]{Loregian}. The universal property of
Definition~\ref{def:lan} follows because maps out of the copower/coend transpose,
via the co-Yoneda density argument, to natural transformations, exactly as in the
proof of Theorem~\ref{thm:density}. \end{proof}

\subsection{Free coproduct completion \emph{is} a left Kan extension}

We now single out the instance the book uses. Recall the container category.

\begin{notation}[$\Cont$ and $\Fam(\Set\op)$]
\label{not:cont}
A morphism of containers $(S,P)\to(T,Q)$ is a pair $(u,v)$ with
$u\colon S\to T$ a map of shapes and $v$ a family
$v_{s}\colon Q(u\,s)\to P(s)$ of position maps \emph{going backwards}. This is
exactly the data of a family, indexed by $S$, of objects of $\Set\op$ (namely the
$P(s)$) together with a map of the index sets and, over it, morphisms in
$\Set\op$. Hence
\[
  \Cont \;\iso\; \Fam(\Set\op),
\]
the \emph{free coproduct completion} of $\Set\op$: objects are $S$-indexed
families of objects of $\Set\op$, and $\Fam(\mathcal{A})$ is the universal
category with all (small) coproducts containing $\mathcal{A}$. Under the
extension functor $\ext{-}$, a family $(P(s))_{s\in S}$ maps to the coproduct of
representables $\sum_{s\in S} y^{P(s)}$ (Remark~\ref{rem:containerrep}).
\end{notation}

\begin{theorem}[Free coproduct completion as $\Lan$; Neil's observation (i)]
\label{thm:neil1}
Let $\mathcal{A}$ be a small category and let
$Y\colon\mathcal{A}\to\Fam(\mathcal{A})$ be the ``singleton'' embedding sending
$a$ to the one-element family $(a)$. Then $\Fam(\mathcal{A})$, with $Y$, is the
free cocompletion of $\mathcal{A}$ under coproducts, and consequently:
\begin{enumerate}
\item[\textup{(a)}] For any functor $F\colon\mathcal{A}\to\mathcal{E}$ into a
category $\mathcal{E}$ with small coproducts, the \emph{coproduct-preserving}
extension $\widehat{F}\colon\Fam(\mathcal{A})\to\mathcal{E}$ defined on objects by
\[
  \widehat{F}\bigl((a_{s})_{s\in S}\bigr) \;=\; \coprod_{s\in S} F(a_{s})
\]
is the left Kan extension $\Lan_{Y}F$, and $\widehat{F}\circ Y\iso F$.
\item[\textup{(b)}] In particular $\widehat{F}$ \emph{preserves representables}:
it sends the singleton family $(a)=Y(a)$ to $F(a)$, i.e.\ to the generator
$F(a)$, with no further identifications.
\end{enumerate}
Specialising $\mathcal{A}=\Set\op$, $F = y^{(-)}\colon\Set\op\to[\Set,\Set]$ (the
representable embedding), and $\mathcal{E}=[\Set,\Set]$, the extension
$\widehat{y^{(-)}}\colon\Cont\iso\Fam(\Set\op)\to[\Set,\Set]$ is exactly the
container-extension functor $\ext{-}$, and it is the left Kan extension of the
representable embedding along the singleton embedding
$Y$.\footnote{A size remark, since $\Set\op$ is \emph{not} small and
Theorem~\ref{thm:lancoend} was stated for small $\mathcal C$. Two things rescue
the specialisation, and neither needs a large coend to converge. First, the
objects of $\Fam(\Set\op)=\Cont$ are indexed by \emph{small} shape sets $S$, so
the coproduct $\coprod_{s\in S}F(a_s)$ that the proof lands on is a small
colimit, which exists in $[\Set,\Set]$. Second --- and this is the clean route
--- part~(a) is really a consequence of the \emph{universal property} of the free
coproduct completion (coproduct-preserving functors out of $\Fam(\mathcal A)$
biject with functors out of $\mathcal A$), which holds for large $\mathcal A$
verbatim; the coend computation above is the concrete face of that universal
property, not a hypothesis we depend on.} Thus
\[
  \ext{-}\;=\;\Lan_{Y}\,y^{(-)}, \qquad
  \ext{S,P} \;=\; \coprod_{s\in S} y^{P(s)} .
\]
\end{theorem}

\begin{proof}
We prove (a) and (b); the specialisation is then read off.

\emph{The coend collapses to the coproduct.} Apply
Theorem~\ref{thm:lancoend} with $K=Y\colon\mathcal{A}\to\Fam(\mathcal{A})$. For an
object $d=(a_{s})_{s\in S}$ of $\Fam(\mathcal{A})$,
\[
  (\Lan_{Y}F)(d) \;\iso\; \int^{a\in\mathcal{A}} \Fam(\mathcal{A})\bigl(Y a,\,d\bigr)\cdot F(a).
\]
By the definition of morphisms in the free coproduct completion, a map from a
\emph{singleton} family $Y a=(a)$ into $d=(a_{s})_{s\in S}$ is a choice of one
index $s\in S$ together with a morphism $a\to a_{s}$ in $\mathcal{A}$:
\[
  \Fam(\mathcal{A})\bigl(Ya,\,(a_{s})_{s}\bigr) \;\iso\;
  \coprod_{s\in S} \mathcal{A}(a,\,a_{s}).
\]
(There is no coproduct \emph{to} match on the domain side --- the domain is a
singleton --- which is exactly why singletons are the generators.) Substituting,
\[
  (\Lan_{Y}F)(d)
   \;\iso\; \int^{a} \Bigl(\coprod_{s\in S}\mathcal{A}(a,a_{s})\Bigr)\cdot F(a)
   \;\iso\; \coprod_{s\in S}\;\int^{a} \mathcal{A}(a,a_{s})\cdot F(a),
\]
using that copowers and the coend (both colimits) commute with the coproduct over
$s$. The inner coend is the co-Yoneda / density formula
(Theorem~\ref{thm:density}) evaluated for the functor $F$ at the object $a_{s}$:
\[
  \int^{a}\mathcal{A}(a,a_{s})\cdot F(a)\;\iso\;F(a_{s}).
\]
Hence $(\Lan_{Y}F)(d)\iso\coprod_{s\in S}F(a_{s}) = \widehat{F}(d)$, naturally in
$d$. This proves (a): the coproduct-preserving extension \emph{is} the left Kan
extension, with unit $\eta\colon F\iso\widehat{F}\circ Y$ the identification of
$F(a)$ with the one-term coproduct.

\emph{Preservation of representables (b).} By (a),
$\widehat{F}(Ya)=\coprod_{s\in 1}F(a)=F(a)$, on the nose. When $F=y^{(-)}$ the
generators $F(a)$ are the representables themselves, so $\widehat{y^{(-)}}$ sends
singletons to representables: it preserves representables.

\emph{Identification with $\ext{-}$.} For $\mathcal{A}=\Set\op$ and
$d=(P(s))_{s\in S}$ a container $(S,P)$, part (a) gives
$\widehat{y^{(-)}}(S,P)=\coprod_{s\in S}y^{P(s)}$, whose value at $X$ is
$\sum_{s\in S}\Set(P(s),X)=\ext{S,P}(X)$ by Remark~\ref{rem:containerrep}. So
$\widehat{y^{(-)}}=\ext{-}=\Lan_{Y}y^{(-)}$. \end{proof}

\begin{remark}[Why this is exactly the reusable form the book needs]
\label{rem:reusable}
Theorem~\ref{thm:neil1} says the container-extension functor is not an ad hoc
gadget: it is \emph{forced} as the unique cocontinuous (coproduct-preserving)
extension of the representable embedding. Two consequences the monoidal chapter
uses without reproof: (i) any structure on $\Set\op$ that we wish to transport to
$\Cont$ need only be defined on generators (singletons) and extended by
coproducts, and the result is automatically the Kan extension; (ii) to check that
a functor out of $\Cont$ agrees with $\ext{-}$ it suffices to check agreement on
representables plus preservation of coproducts. These are precisely the two hooks
Theorem~C hangs on.
\end{remark}

\begin{example}[The extension of a two-shape container]
\label{ex:lan}
Take $S=\{s_{0},s_{1}\}$, $P(s_{0})=1$, $P(s_{1})=2=\{0,1\}$. Then
\[
  \ext{S,P}\;=\;y^{1}\ \amalg\ y^{2}\;\iso\;\mathrm{Id}\ \amalg\ (-)^{2},
  \qquad
  \ext{S,P}(X)\;=\;X\ +\ X^{2}.
\]
This is the functor whose algebras are ``a nullary-of-arity-1 and a binary
operation'' --- e.g.\ the signature functor for magmas-with-a-unary-op. The
point of Theorem~\ref{thm:neil1} is that we did \emph{not} compute a coend to get
this: the completion being free, the Kan extension is literally the coproduct of
representables indexed by shapes, read straight off $(S,P)$.
\end{example}

% =====================================================================
\section{Day convolution}
\label{sec:day}
% =====================================================================

Day convolution transports a monoidal structure from a category to its presheaf
category, in the universal way compatible with Yoneda. On $\Cont$ it degenerates
to ``multiply shapes, combine positions,'' with no coend to compute --- and that
degeneration is the reason the next chapter's monoidal structures are as
tractable as they are.

\begin{definition}[Day convolution]
\label{def:day}
Let $(\mathcal{C},\otimes,I)$ be a small monoidal category. The \emph{Day
convolution} of copresheaves $F,G\colon\mathcal{C}\to\Set$ is
\[
  (F\Day G)(c) \;=\; \int^{c_{1},c_{2}\in\mathcal{C}}
     \mathcal{C}(c_{1}\otimes c_{2},\,c)\times F(c_{1})\times G(c_{2}),
\]
with unit the representable $y^{I}=\mathcal{C}(I,-)$.
\end{definition}

\begin{theorem}[Day~\cite{Day1970}; universal property]
\label{thm:dayuniv}
$(\,[\mathcal{C},\Set],\Day,y^{I})$ is a symmetric\footnote{Symmetric when
$\otimes$ is; in general braided/plain monoidal as $\otimes$ is.} monoidal
category. It is cocontinuous in each variable, and the Yoneda
embedding\footnote{The domain is $\mathcal C\op$, not $\mathcal C$: recall
(\S\ref{sec:yoneda}) that $y^{(-)}\colon c\mapsto\mathcal C(c,-)$ is
\emph{covariant out of $\mathcal C\op$}, a morphism $a\to b$ of $\mathcal C\op$
inducing $y^{a}\Rightarrow y^{b}$ by precomposition. The tensor is unchanged on
objects, so strong monoidality reads $y^{a}\Day y^{b}\iso y^{a\otimes b}$ either
way; we keep the variance honest because the container application turns exactly
on it.}
\[
  y^{(-)}\colon(\mathcal{C}\op,\otimes,I)\longrightarrow([\mathcal{C},\Set],\Day,y^{I})
\]
is \emph{strong monoidal}: there are natural isomorphisms
\[
  y^{c_{1}}\Day y^{c_{2}} \;\iso\; y^{\,c_{1}\otimes c_{2}},
  \qquad y^{I}=y^{I}.
\]
Moreover Day convolution is the \emph{unique} monoidal structure on
$[\mathcal{C},\Set]$ that is cocontinuous in each variable and makes $y^{(-)}$
strong monoidal.
\end{theorem}

\begin{proof}[Reference and the one computation we will reuse]
The full statement is Day~\cite{Day1970}; see also
Loregian~\cite[Ch.~3]{Loregian}. We record only the representable computation, as
the degeneration below leans on it. By the coend formula and the co-Yoneda lemma
(density, Theorem~\ref{thm:density}) applied twice,
\[
  (y^{a}\Day y^{b})(c)
  = \int^{c_{1},c_{2}} \mathcal{C}(c_{1}\otimes c_{2},c)\times
      \mathcal{C}(a,c_{1})\times\mathcal{C}(b,c_{2})
  \iso \mathcal{C}(a\otimes b,\,c) = y^{a\otimes b}(c),
\]
the first coend variable collapsing $c_{1}\rightsquigarrow a$ and the second
$c_{2}\rightsquigarrow b$ by density. Cocontinuity in each variable is immediate:
coends and products-with-a-fixed-set are colimits and preserve colimits in $F$
(resp.\ $G$). Uniqueness: any cocontinuous-in-each-variable monoidal structure is
determined by its values on the dense generators (the representables) by
Remark~\ref{rem:densegen}, and strong-monoidality pins those values to
$y^{a}\Day y^{b}\iso y^{a\otimes b}$; two such structures therefore agree
everywhere. \end{proof}

\subsection{Key degeneration: Day on $\Cont$ needs no coend}
\label{sub:degen}

Now we cash Theorem~\ref{thm:neil1} against Theorem~\ref{thm:dayuniv}. The
category carrying our presheaves is $\Cont\iso\Fam(\Set\op)$, the \emph{free}
coproduct completion of $\Set\op$; its objects are coproducts of representables;
Day is cocontinuous in each variable. So Day is determined on representables and
extended by coproducts --- and on representables it is just $\otimes$ of the
representing objects. No coend survives.

\begin{proposition}[Degenerate Day on containers]
\label{prop:degen}
Let $(\Set,\star,I)$ be any monoidal structure on $\Set$ (below: $\star\in\{+,
\times,\vee\}$). Taking it as the base monoidal category $\mathcal C=(\Set,\star)$
of Definition~\ref{def:day} --- the opposite enters only through the
\emph{embedding} $y^{(-)}\colon\Set\op\to[\Set,\Set]$ and the identification
$\Cont\iso\Fam(\Set\op)$, not through $\star$ --- Day convolution induces a
monoidal structure $\Dayc_{\star}$ on $\Cont$ given \emph{without any coend} by
\[
  p\ \Dayc_{\star}\ q \;=\;\Bigl(\,S_{p}\times S_{q},\ \ (s,t)\mapsto p[s]\star q[t]\,\Bigr),
\]
where $p=(S_{p},p[-])$ and $q=(S_{q},q[-])$ are containers (writing $p[s]$ for the
position set at shape $s$). On extensions this is the Day convolution induced by
$\star$; on representables it restricts to $y^{a}\Dayc_{\star}y^{b}=y^{a\star b}$.
\end{proposition}

\begin{proof}
Write every container as a coproduct of representables,
$p=\coprod_{s\in S_{p}}y^{p[s]}$, $q=\coprod_{t\in S_{q}}y^{q[t]}$
(Theorem~\ref{thm:neil1}). Day convolution is cocontinuous in each variable
(Theorem~\ref{thm:dayuniv}), so it commutes with these coproducts:
\[
  p\Day_{\star}q
  \;\iso\; \Bigl(\coprod_{s}y^{p[s]}\Bigr)\Day_{\star}\Bigl(\coprod_{t}y^{q[t]}\Bigr)
  \;\iso\; \coprod_{s,t}\bigl(y^{p[s]}\Day_{\star}y^{q[t]}\bigr).
\]
By strong-monoidality on representables (the representable computation in the
proof of Theorem~\ref{thm:dayuniv}, with base $\mathcal{C}=(\Set,\star)$ and the
$\Set\op$-valued embedding $y^{(-)}$),
$y^{p[s]}\Day_{\star}y^{q[t]}\iso y^{\,p[s]\star q[t]}$. Hence
\[
  p\Day_{\star}q \;\iso\; \coprod_{(s,t)\in S_{p}\times S_{q}} y^{\,p[s]\star q[t]},
\]
which is exactly the container $\bigl(S_{p}\times S_{q},\,(s,t)\mapsto p[s]\star
q[t]\bigr)$ read off by Theorem~\ref{thm:neil1}. The coend that defined
$\Day_{\star}$ has been discharged entirely by density on the generators plus
cocontinuity; nothing integral remains. \end{proof}

\begin{remark}[Where the coend went]
\label{rem:coendgone}
The coend in Definition~\ref{def:day} is a colimit that identifies ways of
factoring an object as $c_{1}\otimes c_{2}$. On a \emph{free} coproduct
completion the only objects that need testing are the generators (singletons),
and a generator factors through $\otimes$ in exactly one way up to the coend's
identifications --- that is the content of the representable collapse
$y^{a}\Day y^{b}\iso y^{a\otimes b}$. So the coend does no work beyond the
generators, and cocontinuity carries the answer to all coproducts. This is the
same phenomenon as Theorem~\ref{thm:neil1}: ``free completion $\Rightarrow$
structure is what it does on generators.''
\end{remark}

\begin{example}[Three instances side by side]
\label{ex:three}
Take $\star\in\{+,\times,\vee\}$ where $\vee$ is a fixed monoidal structure on
$\Set$ (for concreteness one may read $\vee$ as a chosen ``max-of-arities''
combination; we do \emph{not} classify these here --- that is the business of the
next chapter). Proposition~\ref{prop:degen} gives, for containers
$p=(S_{p},p[-])$, $q=(S_{q},q[-])$:
\[
\renewcommand{\arraystretch}{1.5}
\begin{array}{r|l|l}
  \star & p\ \Dayc_{\star}\ q \text{ (positions at shape }(s,t)) & \text{one-line reading}\\\hline
  + & p[s] + q[t] & \text{a position on the left \emph{or} on the right}\\
  \times & p[s]\times q[t] & \text{a position on the left \emph{and} on the right}\\
  \vee & p[s]\vee q[t] & \text{the chosen combined arity}
\end{array}
\]
In all three cases shapes multiply, $S_{p}\times S_{q}$; only the position-combining
law changes. Concretely with $p=y^{2}$ (one shape, two positions) and $q=y^{3}$:
$p\Dayc_{+}q=y^{5}$, $p\Dayc_{\times}q=y^{6}$, and $p\Dayc_{\vee}q=y^{\,2\vee 3}$.
We stress that we have proved \emph{the formula}, not any classification of which
monoidal structures on $\Cont$ arise this way; the classification (Theorem~C and
its companions) is the next chapter's work, and the present degeneration lemma is
the tool it stands on.
\end{example}

% =====================================================================
\section{Yoneda, monoidal, and closed}
\label{sec:closed}
% =====================================================================

The last piece is closed structure. A closed monoidal category has internal homs
$[-,-]$ right adjoint to the tensor. Neil's observation~(ii) is that closed
structure on a Day category is \emph{forced on representables and extended by
density}: because representables are a dense generator, the internal hom is pinned
down by its behaviour against $y^{c}$ and has no freedom elsewhere. We prove
exactly that, and we are deliberately careful about a temptation it invites and
does not license. It is tempting to conclude that \emph{any} cocontinuous
strong-monoidal functor out of $[\mathcal C,\Set]$ automatically preserves the
internal hom --- ``it agrees on generators, so it agrees everywhere.'' That
inference is \emph{false}: the internal hom is built from a \emph{limit}
(right adjoint), and a cocontinuous functor is under no obligation to preserve
limits. We give a one-line counterexample (Remark~\ref{rem:notfree}). Preservation
is a genuine, additional fact; it holds for the container-extension functor, and
the monoidal chapter establishes \emph{that} directly, using this section only for
the density principle and the representable formula.

\begin{definition}[Closed monoidal category]
\label{def:closed}
A monoidal category $(\mathcal{V},\otimes,I)$ is \emph{closed} if for every object
$b$ the functor $-\otimes b$ has a right adjoint $[b,-]$, the \emph{internal hom},
i.e.\ there are natural isomorphisms
$\mathcal{V}(a\otimes b,\,c)\iso\mathcal{V}(a,\,[b,c])$.
\end{definition}

\begin{theorem}[Closure is determined on representables; Neil's observation (ii)]
\label{thm:neil2}
Let $(\mathcal{C},\otimes,I)$ be a small monoidal category and equip
$[\mathcal{C},\Set]$ with Day convolution $(\Day,y^{I})$
(Theorem~\ref{thm:dayuniv}).
\begin{enumerate}
\item[\textup{(a)}] \emph{(Day is closed.)} $[\mathcal{C},\Set]$ is closed for
$\Day$: each $-\Day G$ has a right adjoint $[G,-]_{\Day}$, and
\[
  [G,H]_{\Day}(c)\;\iso\;\Nat\!\bigl(y^{c}\Day G,\ H\bigr).
\]
The internal hom is thus \emph{determined by its behaviour on representable
first arguments}; on a representable it is a shift,
$[y^{c},H]_{\Day}\iso H(c\otimes-)$, and it is extended to all $G$ by density
(Theorem~\ref{thm:density}).
\item[\textup{(b)}] \emph{(The hom of two representables.)} Evaluated on a pair of
representables,
\[
  [\,y^{a},\,y^{b}\,]_{\Day}(d)\;\iso\;\Nat\!\bigl(y^{d\otimes a},\,y^{b}\bigr)
  \;\iso\;\mathcal C(b,\ d\otimes a).
\]
This is in general a \emph{coproduct} of representables --- a genuine container ---
\emph{not} a single representable. For the Dirichlet tensor
$\star=\times$ on $\Set$ (Proposition~\ref{prop:degen}) it is the container whose
shapes are the maps $b\to a$, i.e.\ the morphisms $\Cont(y^{a},y^{b})$, with $b$
positions apiece; that is precisely the ``shapes are morphisms'' internal hom of
the monoidal chapter.
\end{enumerate}
Part~(a) is Day's~\cite{Day1970}. What part~(a) does \emph{not} assert is that an
arbitrary functor sharing these values preserves the internal hom; see
Remark~\ref{rem:notfree} and Corollary~\ref{cor:contclosed} for the honest form of
the transport the book uses.
\end{theorem}

\begin{proof}
\emph{(a) The Day-closed structure.} Fix $G$. The functor $-\Day G$ is
cocontinuous in its first variable (Theorem~\ref{thm:dayuniv}), and
$[\mathcal{C},\Set]$ is a presheaf category, hence cocomplete and locally
presentable; a cocontinuous endofunctor of a locally presentable category has a
right adjoint (adjoint functor theorem for locally presentable categories,
Riehl~\cite[\S E]{Riehl}). Call it $[G,-]_{\Day}$. To identify it on
representables, use Yoneda and the adjunction:
\[
  [G,H]_{\Day}(c)
  \;\iso\;\Nat\!\bigl(y^{c},\,[G,H]_{\Day}\bigr)
  \;\iso\;\Nat\!\bigl(y^{c}\Day G,\ H\bigr),
\]
the first step by the Yoneda lemma (Theorem~\ref{thm:yoneda}) and the second by
the $(-\Day G)\dashv[G,-]_{\Day}$ adjunction. This proves the displayed formula
and shows $[G,-]_{\Day}$ is \emph{determined} by its behaviour against
representables --- exactly ``determined on representables and extended by
density.'' Taking $G=y^{c}$ representable and reading the same formula in the
first slot,
\[
  [y^{c},H]_{\Day}(d)\iso\Nat\!\bigl(y^{d}\Day y^{c},\,H\bigr)
  \iso\Nat\!\bigl(y^{d\otimes c},\,H\bigr)\iso H(d\otimes c),
\]
by Yoneda and $y^{d}\Day y^{c}\iso y^{d\otimes c}$
(Theorem~\ref{thm:dayuniv}); so $[y^{c},H]_{\Day}\iso H((-)\otimes c)$ is a shift,
as claimed.

\emph{(b) The hom of two representables.} Put $H=y^{b}$ in the shift formula from
part~(a):
\[
  [y^{a},y^{b}]_{\Day}(d)\iso y^{b}(d\otimes a)
  =\mathcal C(b,\,d\otimes a),
\]
using $[y^{a},y^{b}]_{\Day}(d)\iso\Nat(y^{d}\Day y^{a},y^{b})\iso\Nat(y^{d\otimes
a},y^{b})\iso\mathcal C(b,d\otimes a)$ (Yoneda and strong monoidality). Whether
this is again representable is exactly the question of whether $d\mapsto\mathcal
C(b,d\otimes a)$ is a single representable, which it need not be; it is, however,
always a coproduct of representables by density (Theorem~\ref{thm:density}), hence
a container. For $\mathcal C=(\Set,\times)$,
\[
  [y^{a},y^{b}]_{\Day}(d)=\Set(b,\,d\times a)
  =\Set(b,d)\times\Set(b,a)
  =\coprod_{\phi\in\Set(b,a)}\Set(b,d)
  =\coprod_{\Cont(y^{a},y^{b})} y^{b}(d),
\]
using $\Cont(y^{a},y^{b})\iso\Set(b,a)$ (Corollary~\ref{cor:ff}). So the internal
hom is the container with shape set $\Cont(y^{a},y^{b})$ and $b$ positions at each
shape --- the ``shapes are morphisms'' formula the monoidal chapter certifies
(after Niu--Spivak, Exercise~4.78). It is emphatically not $y^{b^{a}}$.
\end{proof}

\begin{remark}[Density determines; it does \emph{not} transport for free]
\label{rem:notfree}
It is tempting to upgrade Theorem~\ref{thm:neil2} to: \emph{any} cocontinuous
strong-monoidal $L\colon[\mathcal C,\Set]\to\mathcal V$ (with $\mathcal V$ closed)
preserves the internal hom, since it agrees with the target's internal hom on the
dense generators. \textbf{This is false}, and the reason is worth internalising,
because it recurs. The would-be argument reduces $G$ to a coproduct of
representables $G=\coprod_{s}y^{c_{s}}$ and uses that internal-hom-out-of-$G$ is a
\emph{product} $\prod_{s}[y^{c_{s}},-]$; but a cocontinuous $L$ preserves
\emph{co}products, and is under no obligation to preserve that product. Internal
hom is a right adjoint --- a limit-shaped thing --- and colimit-preservation buys
nothing against it.

A one-line counterexample. Take $\mathcal C=\mathbf 1$, so
$[\mathcal C,\Set]=\Set$ with Day tensor the cartesian product and internal hom
$[G,H]=H^{G}$. Let $\mathcal V=\mathrm{Vect}_{k}$ (closed under $\otimes_{k}$) and
$L=k[-]$ the free-vector-space functor: it is cocontinuous (a left adjoint) and
strong monoidal, $k[A\times B]\iso k[A]\otimes k[B]$, $k[1]=k$. The comparison is
\[
  \lambda_{G,H}\colon\ k[H^{G}]\ \longrightarrow\ \Hom_{k}\bigl(k[G],k[H]\bigr),
  \qquad (f\colon G\to H)\ \longmapsto\ \bigl(e_{g}\mapsto e_{f(g)}\bigr).
\]
For $G$ infinite and $|H|\ge2$ (say $G=\mathbb N$, $H=\{0,1\}$) it is \emph{not}
surjective: a general linear map $k[\mathbb N]\to k[H]$ assigns arbitrary,
independent images to the countably many basis vectors, whereas a
\emph{finite} linear combination of the ``function'' maps $e_{g}\mapsto e_{f(g)}$
realises only finitely many distinct column-patterns. So $\lambda$ is not
invertible: $L=k[-]$ is strong monoidal but not strong closed. What distinguishes
it from the container-extension functor is that $k[-]$ is not fully faithful,
whereas $\ext{-}$ is; and it is the honest, checkable preservation --- not an
abstract entailment --- that the next corollary rests on.
\end{remark}

\begin{corollary}[The container reading]
\label{cor:contclosed}
For the Dirichlet tensor $\Dayc_{\times}$ on $\Cont\iso\Fam(\Set\op)$,
Theorem~\ref{thm:neil2} pins the internal hom down with no coend: it is determined
on representables by the shift $[y^{a},H]_{\Day}\iso H((-)\times a)$, and on two
representables it is the container
\[
  [y^{a},y^{b}]_{\Day}\;=\;\Bigl(\Cont(y^{a},y^{b}),\ \text{$b$ positions at each shape}\Bigr),
\]
extended to all containers by density. That $\Dayc_{\times}$ is genuinely closed
--- that this candidate internal hom really is right adjoint to $-\Dayc_{\times}q$
--- is the content the monoidal chapter certifies directly (after Niu--Spivak,
Exercise~4.78; machine-checked in \texttt{DirichletClosed.lean}). Chapter~0
supplies only the two pieces that make the formula \emph{forced}: the density
principle and the representable computation above. It does \emph{not} supply
closure ``for free'' from cocontinuity --- Remark~\ref{rem:notfree} shows why no
such free lunch exists. \emph{(This is the honest form the monoidal chapter
cites.)}
\end{corollary}

\begin{example}[The Dirichlet hom of two representables, computed]
\label{ex:ccc}
Take $\star=\times$. A natural first guess for $[y^{a},y^{b}]_{\Day}$ is the
``pointwise exponential'' $y^{\,b^{a}}$ --- and it is \emph{wrong}, instructively
so. Part~(b) computes the honest answer:
\[
  [y^{a},y^{b}]_{\Day}(d)=\Set(b,d\times a)=\underbrace{\Set(b,d)}_{=\,y^{b}(d)}\times\,\Set(b,a)
  =\coprod_{\phi\colon b\to a} y^{b}(d).
\]
So the internal hom is $\coprod_{\Set(b,a)}y^{b}$: a container with $|\Set(b,a)|$
shapes --- one per morphism $y^{a}\to y^{b}$ --- each carrying $b$ positions. Take
$a=2$, $b=1$, $d$ arbitrary: the guess $y^{b^{a}}=y^{1^{2}}=y^{1}=\mathrm{Id}$
gives $[\,\cdot\,](d)=d$, whereas the truth $\Set(1,d\times2)=d\times2$ gives
$2d$; already off by a factor of two. The exponent lives in the \emph{shapes}
(the morphisms), not in the positions --- which is exactly why the monoidal
chapter's internal hom has ``shapes are morphisms.''  We flag --- for the next
chapter --- that at the level of the full functor category $[\Set,\Set]$ this
Dirichlet-closed structure is \emph{not} the container-composition structure
$\tri$; the present chapter asserts only that, whatever Day tensor we transported,
its internal hom is \emph{forced} on representables and extended by density.
\end{example}

% =====================================================================
\section*{What this chapter bought us}
\addcontentsline{toc}{section}{What this chapter bought us}
% =====================================================================

Three reusable facts, each proved from the same two ideas (Yoneda + density):
\begin{itemize}
\item \textbf{Representables generate.} Every functor is a canonical colimit of
representables (Theorem~\ref{thm:density}); representables are a dense generator
of $[\Set,\Set]$.
\item \textbf{Completion is Kan.} The container-extension functor is the left Kan
extension of the representable embedding along the singleton inclusion, i.e.\
free coproduct completion, and it preserves representables
(Theorem~\ref{thm:neil1}). Hence Day convolution on $\Cont$ needs no coend
(Proposition~\ref{prop:degen}): $p\Dayc_{\star}q=(S_{p}\times S_{q},\,(s,t)\mapsto
p[s]\star q[t])$.
\item \textbf{Closure is determined, not free.} A Day tensor is closed, and its
internal hom is \emph{forced} on representables --- $[y^{a},H]_{\Day}\iso
H((-)\otimes a)$, and $[y^{a},y^{b}]_{\Day}$ the ``shapes are morphisms''
container --- then extended by density (Theorem~\ref{thm:neil2}). Density
\emph{determines} the hom; it does not hand it over to every cocontinuous functor
for free (Remark~\ref{rem:notfree}), and the actual closure of the container
tensors is certified directly in the monoidal chapter.
\end{itemize}
These are the pre-container tools the monoidal chapter's Theorem~C and its
companions stand on. Nothing here is new; the arrangement --- and the care about
what density does and does not buy --- is ours.

\begin{thebibliography}{9}
\bibitem{Day1970} B.~Day, \emph{On closed categories of functors}, in Reports of
the Midwest Category Seminar IV, Lecture Notes in Mathematics 137, Springer,
1970, pp.~1--38.
\bibitem{Loregian} F.~Loregian, \emph{(Co)end Calculus}, London Mathematical
Society Lecture Note Series 468, Cambridge University Press, 2021.
\bibitem{MacLane} S.~Mac Lane, \emph{Categories for the Working Mathematician},
2nd ed., Graduate Texts in Mathematics 5, Springer, 1998 (esp.\ Ch.~III on Yoneda
and Ch.~X on Kan extensions).
\bibitem{Riehl} E.~Riehl, \emph{Category Theory in Context}, Aurora: Dover Modern
Math Originals, Dover, 2016.
\end{thebibliography}

\end{document}
